% Two-page briefing note for the supervisor meeting. Not part of report.tex.
% Every number here is copied from report.tex, which verify_report.py checks
% against the run artifacts. If a number changes there, change it here too.
% Build:  cd report && pdflatex brief.tex
\documentclass[11pt,a4paper]{article}
\usepackage[margin=0.8in]{geometry}
\usepackage{booktabs}
\usepackage{parskip}
\usepackage[hidelinks]{hyperref}
\setlength{\parskip}{4pt}
\pagestyle{empty}

\begin{document}

{\large\textbf{AE496 UGP brief: why a physics-informed network loses the cylinder wake, and what brings it back}}\\[2pt]
Peeyush Agarwal \hfill \today

\textbf{The question.} If you train a physics-informed neural network (PINN) on
the flow behind a cylinder without giving it any data, it tends to lose the
vortex street and drift towards the steady, symmetric flow. That steady flow
satisfies the Navier-Stokes equations and every boundary condition exactly, so
the loss has no reason to reject it. Chuang and Barba documented this at
Re = 200. I wanted to reproduce the failure on ground truth I had generated and
validated myself, and then find out which fix actually defeats it, with
controls tight enough that the credit goes to the right mechanism. A second
thread follows the reference paper by Renn et al., whose Fourier Neural
Operator (FNO) is purely data-driven: does adding physics to an operator
forecast help?

\textbf{Ground truth.} I wrote a D2Q9 two-relaxation-time lattice Boltzmann
solver in JAX. At Re = 100 it gives St = 0.1648 against Qu et al.'s 0.1649
(within 0.04\%), with mean drag 3.5\% high and r.m.s.\ lift 6.2\% high. Getting
there took two fixes that I did not expect. First, a domain with free-slip side
walls is an acoustic duct, and given enough time the wake locked onto its
transverse mode (St fell to 0.1132 over 38 periods). A viscous sponge did
nothing; a layer that relaxes the flow towards the free stream removed the
problem. Second, moving the inlet from 8D to 25D upstream cut the drag error
from 7.4\% to 2.3\% (at D = 32), while the domain height barely mattered. An
earlier draft blamed blockage, which was wrong. I then regenerated data at
Re = 60, 80, 120, 150 and 180, and every dataset's St is within 0.6\% of Qu et al.

\textbf{The operator.} The underlying framework, \texttt{underPINN}, only had a
steady residual, so I wrote the unsteady Navier-Stokes operator with a faster
JVP-based derivative path. It annihilates the Taylor-Green vortex to
$1.6\times10^{-16}$, matches the framework's own steady operator to
$5.6\times10^{-17}$, and runs about 2.5 times faster than the literal Hessian.

\textbf{The main result.} I score each run by \emph{shedding retained}: the
r.m.s.\ transverse velocity at a wake probe, PINN over reference. A plain
network trained over one shedding period keeps 31.1\% on average, but anywhere
from 19\% to 43\% depending on where in the cycle the run starts. That phase
dependence turned out to matter a great deal, so every headline number below is
averaged over four start phases with five seeds each. I tested three mechanisms
as switches on one trainer: Fourier features (against spectral bias), a
Reynolds decomposition around a fitted mean flow, and causal time-marching
over four quarter-period windows, where each window starts from the previous
window's prediction so the run sees ground truth only once.

\begin{center}\small
\begin{tabular}{lccc}
\toprule
configuration (data-free, Re = 100) & retained, phase mean & range over phases & velocity error \\
\midrule
plain network, 20\,000 epochs        & 31.1\% & 19 to 43\% & 0.166 \\
plain network, 80\,000 epochs        & 48.6\% & 34 to 63\% & 0.115 \\
time-marching alone                  & 56.4\% & 40 to 74\% & 0.156 \\
all three mechanisms                 & \textbf{72.9\%} & \textbf{70 to 77\%} & \textbf{0.031} \\
\bottomrule
\end{tabular}
\end{center}

The honest reading is more modest than my earlier one. Marching alone gains
only 7.7 points over a plain network given the same 80\,000 epochs, less than
the 17.5 points that longer training adds by itself, and it keeps the phase
swing. Fourier features and the decomposition do not improve retention on
their own. Put all three together, though, and the result stops depending on
start phase, the velocity error drops to a fifth, and marching is worth 34.1
points at equal budget. Over five seeds the phase means carry 95\% intervals of
1.5 points or less for the plain and marched networks and 4.6 for all three. The mechanisms are complementary, which agrees with
Shah (2026). A little data goes a long way: 100 supervised points per window
take retention to 87.4\% and 500 points to 96.9\%.

\textbf{Does it forecast?} Over one period the best data-free setup tracks the
reference phase to 0.1 degrees at 73.5\% amplitude. Over five chained periods it
keeps shedding, but 10.4\% too slowly, and it ends roughly half a cycle out of
phase with 63.8\% of the amplitude. So it sustains a wake; it does not predict
one. Holding the phase takes data, but little: re-anchoring the chain with 350
supervised points per shedding period keeps it within 0.003 of a period of the
reference over all five periods, on both seeds, from 1\,750 points, about 2\% of
the 80\,000 that a fresh snapshot at every window used. At 300 points per period
it still drifts, so the threshold lies between the two. Across Re = 60 to 180 the pattern holds: all three mechanisms keep 63 to
79\% of the wake against 24 to 45\% for the plain network, and cut the velocity
error from about 0.17 to between 0.028 and 0.040.

\textbf{Other architectures.} A GRU-based sequence PINN keeps less of the wake
than the plain network, both on one window (21.8\%) and marched (33.6\%). A single
marched PINNsFormer run reached 77.5\%, but at about twelve times the cost, and
one run at one phase is a lead rather than a result.

\textbf{Physics in operators.} On the validated data the data-only FNO already
forecasts to an error of 0.0010 at its trained horizon, and adding physics makes
it 16 to 18\% worse; the test states are close to ones it has seen, so this is
an easy case. On the earlier, acoustically contaminated dataset physics helps the
FNO by 21\% at the trained horizon but crosses over within about one period. It
leaves a U-Net within 7\% and makes a convolutional GRU 30\% worse, and both
convolutional networks beat the FNO without physics while using 31 and 131
times fewer parameters. Whether physics helps depends on the backbone, not on
operator learning in general.

\textbf{What I got wrong earlier, and fixed.} My first headline claimed marching
alone kept 89.2\% of the wake. Two things inflated it: the marched runs were
given a fresh reference snapshot at every window (worth about twenty points),
and every configuration was scored at a single start phase. A five-period run
that ``did not drift'' was being re-anchored to the truth at every window, and
the ground truth itself was being pulled by the duct mode. Correcting these
reversed the conclusion that one mechanism is enough.

\textbf{Limitations and next steps.} The headline has five seeds, but the
sweeps, the architecture study and the re-anchoring runs have two; the
equal-budget control uses a single budget; the Re sweep uses one start phase
per Re; the U-Net and ConvGRU were only run on the earlier dataset; drag stays
3.5 to 5.7\% high across the sweep (nothing in the results depends on it).
Whether the re-anchoring threshold moves with a longer horizon or another
Reynolds number is untested. Enforcing window continuity exactly (Roy and
Castonguay) would remove the small kinks at window handoffs. The full
PINNsFormer grid would take about 66 GPU-hours.

\textbf{Verification.} Every number in the full report is registered against
the run that produced it, and \texttt{verify\_report.py} currently reproduces
all 557 of them from the saved artifacts.

\end{document}
